\documentclass[11pt]{article}
\usepackage[T1]{fontenc}
\usepackage{microtype}
\usepackage{lmodern,amsmath,amssymb,amsthm,geometry,hyperref}
\geometry{margin=1in}
\hypersetup{hidelinks,pdftitle={Validated Dirichlet approximation and weighted arithmetic moments},pdfauthor={GRH research working note}}
\newtheorem{theorem}{Theorem}
\newtheorem{proposition}[theorem]{Proposition}
\title{Validated Dirichlet approximation and weighted arithmetic moments}
\author{GRH research working note}
\date{5 November 2026}
\begin{document}\maketitle
\begin{abstract}Exact outward dyadic arithmetic and theta-integral quadrature certify three finite Dirichlet readouts. Separately, a fixed-twist sextic Hecke-family moment improves an original masked remainder estimate, while retaining a positive cutoff loss. We give the precise external analytic inputs and distinguish these results from zero-location and spectral determinant claims. No proof of GRH or complete Lean formalization is asserted.\end{abstract}
\section{A finite, validated approximation result}
Let $\chi_3=(0,1,-1)$, $\chi_4=(0,1,0,-1)$ and
$\chi_5=(0,1,-1,-1,1)$, extended periodically, with parities $a=1,1,0$.
We approximate the completed functions
\[
 \Lambda_q(s)=(q/\pi)^{(s+a)/2}\Gamma((s+a)/2)L(s,\chi_q).
\]
The primitive theta/Mellin identity is the explicit classical analytic input;
no hypothesis about zeros is used. These three examples do not assert a
uniform result for all conductors, general Hecke functions or Artin holomorphy.

\begin{theorem}[Scope of the supplied finite certificate]
For each $q=3,4,5$, the supplied exact-rational calculation constructs a
complex polynomial $P_q$ of degree $28$ satisfying
\[
 \sup_{|s-1/2|\le7/10}|P_q(s)-\Lambda_q(s)|<5\cdot10^{-15}.
\]
The computation includes enclosures of all $64$ true Fourier-circle nodes,
completed-function sample error, analytic truncation, Fourier aliasing,
the exact closed-form $180$-step gradient-descent iterate and coefficient
midpoint rounding. The Python implementation and informal analytic proof
have independent checks; this is not a kernel-checked Lean theorem.
\end{theorem}
\begin{proof}
Write $F(c+z)=\sum_{n\ge0}a_nz^n$, with $c=1/2$, $R=7/10$,
$\rho=6/5$, $r=12/5$. The independently recomputed rational theta majorants
are strictly below $M_q=4,10,6$, respectively, on $|z|\le r$.
Use true roots $\zeta_j=e^{2\pi i j/64}$ and rational labels $y_j$
with proved error $|y_j-F(c+\rho\zeta_j)|\le\eta$.
Use the real Euclidean gradient on complex coefficient coordinates for
\[
 \mathcal L(b)=\frac1{128}\sum_{j=0}^{63}
 \left|\sum_{n=0}^{28}b_n\zeta_j^n-y_j\right|^2.
\]
Discrete Fourier orthogonality gives the least-squares optimum
$b_n^*=64^{-1}\sum_jy_j\bar\zeta_j^n$, $0\le n\le28$,
and gradient $\nabla\mathcal L(b)=b-b^*$.
For step $1/4$ and zero initialization, the exact $180$-step iterate is
$b_n^{180}=(1-(3/4)^{180})b_n^*$. The code evaluates this closed form;
it does not claim an explicit180-update optimizer run.

Set $u=R/\rho=7/12$, $v=R/r=7/24$ and $w=(\rho/r)^{64}=2^{-64}$.
Let $\widehat b_n$ be an outward rectangle enclosing $b_n^*$ and
$\epsilon_n$ the sum of the coordinate halfwidths of the rectangle
for $b_n^{180}$. Choose each actual polynomial coefficient as that
rectangle's exact rational midpoint $c_n$. The normalized basis is explicit:
\[
 P_q(s)=\sum_{n=0}^{28}c_n\left(\frac{s-1/2}{6/5}\right)^n.
\]
Thus the exported coefficient with index $n$ multiplies the normalized
variable, rather than the raw power $(s-1/2)^n$. Its error on the compact disk is bounded by
\begin{align*}
 E_{\rm opt}&=(3/4)^{180}\sum_{n=0}^{28}|\widehat b_n|_1u^n,&
 E_{\rm arithmetic}&=\sum_{n=0}^{28}\epsilon_nu^n,\\
 E_{\rm oracle}&=\eta\frac{1-u^{29}}{1-u},&
 E_{\rm alias}&=M_q\frac{w}{1-w}\frac{1-v^{29}}{1-v},\\
 E_{\rm Taylor}&=M_q\frac{v^{29}}{1-v}.&&
\end{align*}
Here $|\widehat b_n|_1$ is the sum of the maximum absolute real and imaginary
coordinates, a safe upper bound for every enclosed complex modulus.
Cauchy's estimate $|a_n|\le M_qr^{-n}$ and the Fourier identity
$b_n^*=a_n\rho^n+\sum_{\ell\ge1}a_{n+64\ell}\rho^{n+64\ell}$
for exact labels prove the alias and Taylor bounds. Label perturbations
contribute at most $\eta$ per Fourier coefficient. The remaining two
terms follow directly from the exact iterate and coefficient rectangles.
Each term is rounded upward to a $180$-bit dyadic rational; the reported total
is the exact sum of those reported terms. The certified totals below are
strictly below the asserted threshold. Detailed sample enclosures are proved
in the extended edition and encoded with exact integer endpoints.
\end{proof}

\begin{center}\begin{tabular}{c|cc}Conductor & Maximum sample radius & Uniform polynomial bound\\\hline
3 & $3.026\times10^{-17}$ & $1.786\times10^{-15}$ \\
4 & $7.206\times10^{-17}$ & $4.455\times10^{-15}$ \\
5 & $1.404\times10^{-16}$ & $2.906\times10^{-15}$ \\
\end{tabular}\end{center}
The displayed values are rounded upward; exact rational records establish the strict inequalities.


\section{What this does and does not certify}
The final JSON records exact rational centers, interval endpoints, source
hashes and all five error terms. Float values are display summaries only.
A high-precision comparison is an independent diagnostic, not the reason
an interval contains its mathematical target. Arithmetic uses outward
integer operations, and invalid float/order inputs are rejected, including
through typed caches. Earlier pilots and failed validation probes are retained.

This is a fixed finite tolerance/configuration. Increasing accuracy needs
new truncation and precision choices; the fixed $N=11,T=128$ interface is
not an arbitrary-accuracy oracle. The broader approximation existence theorem
is separate. The certificate says nothing about the location of zeros:
$P_q$ is not required to be real-rooted or a self-adjoint spectral determinant.
The independent arithmetic determinant identity, positivity and original
signed short-row/covariance estimates remain open. No GRH proof or complete
new Lean4 formalization is claimed.

\section{Requested-error readouts from preserved certified labels}\label{ar:requested}
The accepted final-v5 sample certificate is preserved unchanged. For each
$q=3,4,5$, it supplies $m=64$ rational complex labels $y_j$, bounds
$|y_j-\Lambda_q(c+\rho\zeta_j)|\le\eta_j$, and outward rectangles enclosing
the true roots $\zeta_j=e^{2\pi i j/m}$, where $c=1/2$, $\rho=6/5$.
The classical primitive theta/Mellin identity and accepted growth bounds
$M_q=4,10,6$ on $|s-c|\le r=12/5$ remain the analytic inputs.
This extension accepts an exact rational error request $\varepsilon>0$ on
the fixed compact $|s-c|\le R=7/10$. It changes the finite readout, while
retaining the certified sample errors and the $64$-node alias budget.
It is not an arbitrary-accuracy theta-sampling interface.

\begin{theorem}[Executed rational readout with explicit refusal]\label{ar:returned}
For each supported conductor and rational request, the finite-data
selection below either returns a normalized rational polynomial $P$ and an
executed, exactly replayable gradient-descent trace with
\[
 \sup_{|s-c|\le R}|P(s)-\Lambda_q(s)|\le\varepsilon,
\]
or refuses the request with the exact minimum base budget over all
$64$ allowed degrees. Refusal concerns attainability by this particular
fixed-data majorant; it is not an impossibility theorem for approximation.
\end{theorem}
\begin{proof}
Let $B_n$ be an outward rectangle enclosing
$b_n^*=m^{-1}\sum_j y_j\bar\zeta_j^n$, $0\le n<64$, computed from the
preserved root rectangles. Put $t_n=\operatorname{mid}(B_n)$ and let
$e_n$ be the sum of its coordinate halfwidths. These are exact rationals,
and $|t_n-b_n^*|\le e_n$.
For a selected $d<64$ execute, using exact rational operations,
\[
 b_n^0=0,\qquad b_n^{k+1}=\tfrac34 b_n^k+\tfrac14 t_n,
 \qquad 0\le n\le d.
\]
This is genuine gradient descent with step $1/4$ for the explicitly
specified real Euclidean objective
$J(b)=\frac12\sum_{n=0}^d|b_n-t_n|^2$, whose gradient is $b-t$.
All iterates and exact objective values are retained. The output is literally
\[
 P_k(s)=\sum_{n=0}^d b_n^k\left(\frac{s-c}{\rho}\right)^n.
\]
Set $u=R/\rho=7/12$, $v=R/r=7/24$, $w=(\rho/r)^{64}=2^{-64}$,
$S_d(x)=\sum_{n=0}^d x^n$ and $\eta=\max_j\eta_j$. The five budgets are
\begin{align*}
 E_{\rm opt}(k)&=\sum_{n=0}^d|b_n^k-t_n|_1u^n,&
 E_{\rm oracle}&=\eta S_d(u),\\
 E_{\rm alias}&=M_q\frac{w}{1-w}S_d(v),&
 E_{\rm Taylor}&=M_q\frac{v^{d+1}}{1-v},\\
 E_{\rm arithmetic}&=\sum_{n=0}^d e_nu^n.&&
\end{align*}
Here $|x+iy|_1=|x|+|y|$. Cauchy's estimate and true-root Fourier
orthogonality give the alias and Taylor terms; certified label errors give
the oracle term. Triangle inequality applied to
$b_n^k-b_n^*=(b_n^k-t_n)+(t_n-b_n^*)$ gives the other two.
Thus their exact sum bounds the compact-uniform error.

Define $C_d=E_{\rm oracle}+E_{\rm alias}+E_{\rm Taylor}
+E_{\rm arithmetic}$. Select the smallest $d<64$ with $C_d<\varepsilon$;
if none exists, report $\min_{d<64}C_d$ and refuse. For a selected degree,
$E_{\rm opt}(k)=(3/4)^k\sum_n|t_n|_1u^n\longrightarrow0$.
Therefore executing updates until
$E_{\rm opt}(k)\le\varepsilon-C_d$ terminates after finitely many steps.
The serialized coefficients, budgets, targets and trace use exact rational
endpoints; displayed decimal approximations are not proof inputs.
\end{proof}

% Accepted prose AFM proof; explicit external inputs, not a Lean theorem.
\section{An arithmetic-family moment with every original row retained}
\label{mm:core}
Let $K=\mathbb Q(\sqrt{-3})$, $\mathcal O=\mathbb Z[\omega]$, and write
$N$ for the ideal or element norm. Fix a finite-order Hecke character
$\nu$ and a finite set $S$ of prime ideals containing those above $6$ and
those dividing its conductor. For $0\ne u\in\mathcal O$, let $\kappa_u$
be the character of $K(u^{1/6})/K$ defined using arithmetic Frobenius, and
put $\psi_u=\operatorname{prim}(\nu\kappa_u)$. Here
$\operatorname{prim}$ means the underlying primitive Hecke character:
it is formed before introducing any divisor variable or evaluating
conductor-dependent zero extensions. The rows below are all actual
Eisenstein elements, including units and prime powers supported in $S$.

The external inputs are abelian Kummer reciprocity and character
uniqueness, the primitive Hecke functional equation and its usual
vertical growth, Stirling's formula, scalar Dedekind-zeta convexity, and
Blomer--Goldmakher--Louvel's large sieve for their specified sixth-order
family. We use the pinned
\href{https://arxiv.org/pdf/1112.1650v1}{arXiv:1112.1650v1},
Theorem~1.3, Lemma~2.1 and Appendix~II; its Section~5 supplies the
all-ideal decomposition whose fixed-twist proof is redone here.
No statement about all primitive characters of comparable conductor,
Artin holomorphy, Lindel\"of bounds or zero locations is assumed.

\begin{theorem}[The original-family numerator moment]
\label{mm:afm}
For every $\varepsilon>0$, uniformly for $H\ge1$ and $t\in\mathbb R$,
\begin{equation}\label{mm:central}
 \sum_{0<Nu\le H}|L(1/2+it,\psi_u)|^2
 \ll_{\nu,S,\varepsilon}H^{1+\varepsilon}(1+|t|)^{1+\varepsilon}.
\end{equation}
For every fixed $0<\delta<1/2$, the same estimate holds with
$1/2+it$ replaced by $1/2+\delta+it$, with the constant also depending
on $\delta$. If the central bound has constant $C_0$ at an exponent
$0<\varepsilon_0\le1/4$, an admissible off-central constant is
\begin{equation}\label{mm:poisson-constant}
 200 C_0\left(\frac{3/2+\delta}{1/2-\delta}\right)^2.
\end{equation}
\end{theorem}
\begin{proof}
For the canonical BGL character $\vartheta_{\mathfrak a}$,
Appendix~II gives a radicand $\alpha_{\mathfrak a}$ with
$(\alpha_{\mathfrak a})_S=\mathfrak a\mathfrak g^{-6}$.
Since $\mathcal O_S$ is principal, comparison with any literal $u$
whose outside-$S$ ideal is $\mathfrak a$ gives
$\kappa_u=\eta\vartheta_{\mathfrak a}$ for one of finitely many fixed
$S$-unit Kummer twists $\eta$. This is an equality of primitive
characters, established outside a finite set and extended by uniqueness.
For each outside-$S$ ideal $\mathfrak a$ and $S$-supported integral
ideal $\mathfrak b$, at most six elements generate
$\mathfrak a\mathfrak b$. The resulting $S$-supported multiplicity sum
converges at exponent $1+\varepsilon$.

For each fixed $\xi=\nu\eta$, split the canonical parameter as
$\mathfrak a=\mathfrak a_0\mathfrak a_1\mathfrak a_2^6$, where
$\mathfrak a_0$ is squarefree, $\mathfrak a_1$ is squarefull and
sixth-power-free, and $(\mathfrak a_0,\mathfrak a_1)=1$.
The primitive conductor is at most a fixed multiple of
$N(\mathfrak a_0\mathfrak a_1)$.
The common-gamma approximate functional equation, with
$X=H/(N\mathfrak a_2)^6$ and $T=1+|t|$, has common maximal length
$B=X^{1/2}T(HT)^\eta$ for any sufficiently small fixed $\eta>0$.
Mellin inversion removes its row-dependent conductor weight.
Split a column as $\mathfrak b_S\mathfrak b_1\mathfrak b_2$, with
$\mathfrak b_S$ supported in $S$, $\mathfrak b_1$ squarefree and
$\mathfrak b_2$ squarefull, the last two coprime and prime to $S$.
The original primitive values at $\mathfrak b_S$ remain bounded row
multipliers. The $S$-supported and squarefull column weights are summable
at $1/2+\eta$. Multiplying squarefree column coefficients by the fixed
$\xi$ preserves their squared norm, whatever its finite order.
BGL's large sieve therefore bounds the energy for fixed
$\mathfrak a_1,\mathfrak a_2$ by
\[
 (HT)^{O(\eta)}\{M+B+(MB)^{2/3}\},\qquad
 M=X/N\mathfrak a_1.
\]
Squarefull counting and summation give
$X+B\sqrt X+(XB)^{2/3}\ll (HT)^{O(\eta)}XT$.
The $\mathfrak a_2$ sum converges at exponent $6$.
Principal rows are separated before the pole-free functional equation
and bounded by scalar zeta convexity. This proves~\eqref{mm:central}.

For the off-central statement divide each $L(s,\psi)$ by
$(s+1)^{5/4}$ and, on principal rows, also multiply by $(s-1)/(s+1)$.
These are bounded analytic functions on $\operatorname{Re}s>1/2$,
continuous on the boundary. Poisson representation and Jensen's
inequality transfer the central squared moment. The convolution of
$P_\delta(y)=\delta/[\pi(\delta^2+y^2)]$ with
$(1+|t|)^{-3/2+\varepsilon_0}$ is at most eleven times that weight
when $0<\varepsilon_0\le1/4$. Restoring the denominator and the
principal pole-removal factor gives~\eqref{mm:poisson-constant}.
The extended edition gives the detailed derivation with all primitive local values.
\end{proof}

\paragraph{Separate fixed-unitary extension.}
The moment theorem also holds for any one fixed unitary Hecke character
$\nu$. For precision choose
$\nu(\mathfrak a)=\nu_0(\mathfrak a)(N\mathfrak a)^{-i\tau}$, with
angular parameter $m\in\mathbb Z$ and completed gamma factor
\[
 \gamma_\nu(s)=2(2\pi)^{-(s+|m|/2+i\tau)}
                    \Gamma(s+|m|/2+i\tau).
\]
All finite Kummer twists have this same infinity type. Use
$T_\nu=1+|t+\tau|+|m|$, comparable to $1+|t|$ with
$\nu$-dependent constants. Pure norm rows have
$L(s,\psi)=\zeta_K(s+i\tau)$ and are separated before the pole-free
functional equation; their pole at $1-i\tau$ is removed by
$(s-1+i\tau)/(s+1+i\tau)$. No uniformity in $m$, $\tau$ or the
conductor of $\nu$ is asserted. This extension of the moment is separate
from the following finite-order type-I residue formula.

\subsection{The exact masked remainder and the remaining loss}
Retain the original zero-extended symbol $\chi_d(u)=(u/d)_6$ on good
squarefree ideals. Let $\mu$ be the ideal M\"obius function and set
\[
 \beta_U(d)=\sum_{ab=d\atop Na,Nb\le U}\mu(a)\mu(b),\quad
 T(u,d)=S\cup\operatorname{supp}(u)\cup\operatorname{supp}(d).
\]
Let $\psi_{u,2}=\operatorname{prim}(\psi_u^2)$, formed independently
of squaring any zero sentinel. With $\mathfrak f_u$ and
$\mathfrak f_{u,2}$ the primitive conductors, the full-mask series is
\begin{align}
 F_{u,d}(s)
 &=\frac{L(s,\psi_u)}{L(2s,\psi_{u,2})}C_{T(u,d)}(s),\label{mm:masked}\\
 C_T(s)
 &=\frac{\prod_{p\in T,\ p\nmid\mathfrak f_u}
              (1-\psi_u(p)(Np)^{-s})}
         {\prod_{p\in T,\ p\nmid\mathfrak f_{u,2}}
              (1-\psi_{u,2}(p)(Np)^{-2s})}.\notag
\end{align}
Initially it equals
$\sum_{c\ \mathrm{squarefree},\ (c,T)=1}\psi_u(c)(Nc)^{-s}$
for $\operatorname{Re}s>1$. Every prime of $T$ remains deleted,
including those forgotten by either primitive conductor.

For finite-order $\nu$, $D\ge1$, $U\ge1$, $U^2\le D$, and
$W\in C_c^\infty((0,\infty))$ supported in $[1,2]$, put
\[
 \widehat W(s)=\int_0^\infty W(x)x^{s-1}\,dx,\qquad
 \mathcal M_4(W)=\sum_{j=0}^4\sup_x|(x\partial_x)^jW(x)|.
\]
On $\sigma=1/2+\delta$, the exact residue-subtracted type-I remainder is
\begin{equation}\label{mm:remainder}
 E_u=\frac{D^\delta}{2\pi}\int_{\mathbb R}
 \widehat W(\sigma+it)D^{it}
 \sum_{d\ \mathrm{squarefree}\atop(d,S)=1}
 \frac{\beta_U(d)\nu(d)\chi_d(u)}{(Nd)^{\sigma+it}}
 F_{u,d}(\sigma+it)\,dt.
\end{equation}
This definition retains the actual $\beta_U$, all phases, $d$-masks,
conductors and the independently primitive square.

\begin{theorem}[An improved masked residual estimate]
\label{mm:residual}
Under the preceding finite-order hypotheses, for every $\varepsilon>0$,
\begin{equation}\label{mm:energy}
 \sum_{0<Nu\le H}|E_u|^2
 \ll_{\nu,S,\delta,\varepsilon}
 \mathcal M_4(W)^2D^{2\delta}U^{2-4\delta}H(HD)^\varepsilon.
\end{equation}
\end{theorem}
\begin{proof}
The same deleted set in both factors of~\eqref{mm:masked} gives
$|F_{u,d}(\sigma+it)|\ll (Nu\,Nd)^\eta|L(\sigma+it,\psi_u)|$:
the reciprocal incomplete Euler product at $2\sigma>1$ is bounded
by $\zeta_K(2\sigma)$, and its remaining finite numerator is bounded
by a small power of its radical. Apply Minkowski to~\eqref{mm:remainder}
in the exact row space and use Theorem~\ref{mm:afm}. The actual coefficient
sum is at most
\[
 \sum_d\frac{|\beta_U(d)|}{(Nd)^{\sigma-\eta}}
 \le\left(\sum_{Na\le U}(Na)^{-\sigma+\eta}\right)^2
 \ll U^{1-2\delta+2\eta}.
\]
Four Mellin integrations by parts bound the remaining height integral
by a constant times $\mathcal M_4(W)$. Square the result and absorb
the small losses using $U^2\le D$.
\end{proof}

At $H=D^{3/4}$ and $U=D^a$ with fixed $0<a\le1/2$, the upper bound's excess
over $H$ is $2\delta+(2-4\delta)a$, up to small losses; as
$\delta\downarrow0$ it still leaves $2a$. This is a limitation of the
upper bound, not an energy lower bound. The extracted principal residue,
the original short term at smaller scales and the signed type-II minus
type-I covariance are not controlled by~\eqref{mm:energy}. The next
arithmetic task is to estimate the complete $d$-weighted moment before
its absolute coefficient sum, or establish sufficient signed covariance.
No positive arithmetic determinant, zero-location theorem or GRH proof
follows. The analytic argument uses the external classical inputs above;
finite diagnostics and selected existing Lean local identities do not
formalize the Hecke functions, approximate functional equation, Poisson
transfer or these moment bounds.

\section{Recombining the actual weighted cutoff}
\label{rc:core}
Keep the finite-order Eisenstein family and original masks of
Section~\ref{mm:core}. For $U\ge1$ and a squarefree ideal $m$, put
\[
 A_U(m)=\sum_{d\mid m}\beta_U(d)
 =\sum_{ab\mid m\atop(a,b)=1,\ Na,Nb\le U}\mu(a)\mu(b).
\]
Every summation variable here is an integral ideal. The last formula
uses squarefreeness of $m$; an unconstrained square of a M\"obius
polynomial would also include forbidden overlapping supports.
\begin{theorem}[Exact recombination and an absolute-half-plane rate]
\label{rc:recombination}
Write $T_0=S\cup\operatorname{supp}(u)$ and
$J_{u,U}(s)=\sum_{d\ \mathrm{squarefree},\ (d,S)=1}
\beta_U(d)\nu(d)\chi_d(u)(Nd)^{-s}F_{u,d}(s)$.
For $\operatorname{Re}s>1$,
\begin{align*}
 J_{u,U}(s)&=\sum_{m\ \mathrm{squarefree},\ (m,T_0)=1}
                   A_U(m)\psi_u(m)(Nm)^{-s},\\
 |A_U(m)|&\le3^{\omega(m)},\qquad
 A_U(m)=\mu(m)\quad(Nm\le U).
\end{align*}
For every $\kappa>0$ and $0<\epsilon<\kappa$, uniformly for
$\operatorname{Re}s\ge1+\kappa$, all original rows and all $U\ge1$,
\[
 \left|J_{u,U}(s)-\frac1{L^{T_0}(s,\psi_u)}\right|
 \le U^{-\kappa+\epsilon}
 \{\zeta_K(1+\epsilon)^3+\zeta_K(1+\epsilon)\}.
\]
Here $L^{T_0}$ removes every prime of $T_0$, including conductor-drop
primes; the bound is in the absolute-convergence half-plane only.
\end{theorem}
\begin{proof}
On a surviving term $m=dc$, the two squarefree ideals $c,d$ are
coprime, and both avoid $T_0$. The original product
$\nu(d)\chi_d(u)\psi_u(c)$ is then $\psi_u(m)$. Grouping by $m$
gives $A_U$. Each prime of $m$ has three possible positions in
$(a,b,m/(ab))$, so the absolute sum is at most $3^{\omega(m)}$.
When $Nm\le U$ all divisors fit both cutoffs; the primewise sum is
$1-1-1=-1$, giving $\mu(m)$.

The same bound justifies absolute convergence and regrouping for
$\operatorname{Re}s>1$, because on squarefree ideals
$3^{\omega(m)}=\tau_3(m)$ and
$\sum_m\tau_3(m)(Nm)^{-v}=\zeta_K(v)^3$ for $v>1$,
where this latter sum runs over all nonzero integral ideals.
Subtract the absolutely convergent M\"obius series for $1/L^{T_0}$.
Only $Nm>U$ remain. Rankin's inequality with exponent
$\eta=\kappa-\epsilon$ bounds their absolute sum by
\[
 U^{-\eta}\sum_m\{\tau_3(m)+1\}
                  (Nm)^{-1-\epsilon},
\]
which proves the stated constant independently of deleted primes,
row multiplicities or character phases.
\end{proof}

For $W$ supported in $[1,2]$, the same exact regrouping gives the
original normalized type-I polynomial
\[
 I_u=D^{-1/2}\sum_{m\ \mathrm{squarefree},\ (m,S)=1}
 A_U(m)\nu(m)\chi_m(u)W(Nm/D).
\]
This formula retains all original zero sentinels. It permits a direct
family large sieve before any absolute $d$ sum; its quantitative moment
is developed separately below. The displayed absolute-half-plane rate
alone is not that critical-strip moment.

% Original masked cutoff family; no new Lean or GRH theorem.
\section{A cutoff-uniform envelope for the original masked remainder}
\label{wc:core}
Fix $K=\mathbb Q(\omega)$, $\mathcal O=\mathbb Z[\omega]$, a finite-order
Hecke character $\nu$, and a finite set $S$ containing the primes above $6$
and the conductor of $\nu$. All unspecified ideal sums below are over
nonzero integral ideals. Let $\chi_m(u)=(u/m)_6$ be the original symbol,
zero when $(m,u)>1$, and put
\[
 \beta_U(d)=\sum_{ab=d\atop Na,Nb\le U}\mu(a)\mu(b),\qquad
 A_U(m)=\sum_{d\mid m}\beta_U(d).
\]
The columns $d,c,m$ are squarefree and prime to $S$ where indicated.
Retain the primitive characters $\psi_u=\operatorname{prim}(\nu\kappa_u)$,
$\psi_{u,2}=\operatorname{prim}(\psi_u^2)$, and the full deleted set
$T(u,d)=S\cup\operatorname{supp}(u)\cup\operatorname{supp}(d)$.
For $H,D,U\ge1$, $U^2\le D$, and $W\in C_c^\infty((0,\infty))$
supported in $[1,2]$, the original normalized type-I piece is
\begin{equation}\label{wc:polynomial}
 I_u=D^{-1/2}\sum_{m\ \mathrm{squarefree}\atop(m,S)=1}
     A_U(m)\nu(m)\chi_m(u)W(Nm/D).
\end{equation}
Its exact principal residue is
\begin{align}\label{wc:principal}
 P_u&=\mathbf1_{\psi_u=1}\sqrt D\,\widehat W(1)
 \sum_{d\ \mathrm{squarefree}\atop(d,S)=1}
 \frac{\beta_U(d)\nu(d)\chi_d(u)}{Nd}\rho_{T(u,d)},\notag\\
 \rho_T&=\frac{\operatorname*{res}_{s=1}\zeta_K(s)}{\zeta_K(2)}
               \prod_{p\in T}(1+(Np)^{-1})^{-1}.
\end{align}
Define $E_u=I_u-P_u$; this is the same residue-subtracted remainder as the
fully masked Mellin contour, with no cutoff limit or profile supremum.

\begin{theorem}[Uniformity in the actual truncated cutoff]
\label{wc:uniform}
For every $\varepsilon>0$,
\begin{equation}\label{wc:bound}
 \sum_{0<Nu\le H}|E_u|^2
 \ll_{\nu,S,\varepsilon}\|W\|_\infty^2(HD)^\varepsilon
       \{H+DH^{1/3}+(HD)^{2/3}\}.
\end{equation}
The sum is over every actual nonzero Eisenstein element, including units,
$S$-supported powers and sixth-power rows. The constant is independent of
$U$ under $U^2\le D$. This is an application of the classical sextic
large sieve with exact original coefficients, not a zero-location theorem.
\end{theorem}
\begin{proof}
On an absolutely convergent line, expand the restored squarefree series
$F_{u,d}$. Its condition $(c,T(u,d))=1$ includes $(c,d)=1$, and the
original symbols multiply. Grouping $m=cd$ gives~\eqref{wc:polynomial}
by Mellin inversion. Each prime of $m$ may belong to the first truncated
M\"obius divisor, to the second, or to neither. Consequently
$|A_U(m)|\le3^{\omega(m)}\ll_\eta(Nm)^\eta$, uniformly in $U$.
The length-$D$ coefficient sequence in~\eqref{wc:polynomial} therefore
has squared norm $\ll_\eta\|W\|_\infty^2D^\eta$.

The all-ideal row norm in Section~3 of
\href{https://arxiv.org/pdf/1112.1650v1}{Blomer--Goldmakher--Louvel,
arXiv:1112.1650v1} gives
\[
 B_2^1(X,D)\ll_\eta(XD)^\eta
                 \{X+DX^{1/3}+(XD)^{2/3}\}.
\]
For completeness, split $a=a_1a_2^2\cdots a_5^5a_6^6$ dyadically;
$a_1,\ldots,a_5$ are squarefree and pairwise coprime. If their dyadic
sizes are $A_j$, put $R=\prod_{j=1}^6A_j$ and $M=\max(A_1,A_2)$.
Pivot on the first or second variable, using the squarefree powered
sieve proved simultaneously in that paper's Section~3.2. The other
variables are bounded column multipliers. The block energy is at most
\[
 (XD)^\eta\{R+DR/M+D^{2/3}R/M^{1/3}\}\sum_m|\lambda_m|^2.
\]
The inequalities $R\le X$, $R^3\le XM^3$ and $R^3\le X^2M$
follow from $\prod_jA_j^j\le X$ and the pivot choice. Summing the
logarithmically many blocks proves the displayed bound.

Finite $S$-unit and ray-class partitions convert the original literal
symbol into the canonical column value $\vartheta_m(a)$ with fixed
modulus-one coefficient multipliers. On noncoprime pairs both values
are zero: squarefree $m$ has outside-$S$ conductor exactly $m$.
This orientation preserves sixth-power zeros. At most six elements
generate each outside-$S$ ideal times an $S$-supported ideal; the
$S$-supported sums converge at powers $1$, $1/3$ and $2/3$.
Thus~\eqref{wc:bound} holds first with $I_u$ in place of $E_u$.

A principal $\psi_u$ requires every outside-$S$ valuation of $u$ to
be divisible by $6$. There are $O_{\nu,S}(H^{1/6})$ such literal
rows. Equation~\eqref{wc:principal}, $|\rho_T|\ll_K1$, and the actual
coefficient sum
\[
 \sum_d\frac{|\beta_U(d)|}{Nd}
 \le\left(\sum_{Na\le U}\frac1{Na}\right)^2\ll\log^2(2U)
\]
give $\sum|P_u|^2\ll DH^{1/6}\|W\|_\infty^2\log^4(2U)$.
This is absorbed by the right side of~\eqref{wc:bound}, using
$U^2\le D$ and a smaller internal exponent. Finally
$|E_u|^2\le2|I_u|^2+2|P_u|^2$. The extended proof specifies the
canonical-family transfer and the powered-sieve input precisely.
\end{proof}

Combining this with the previously proved fixed-twist AFM transfer yields
\begin{equation}\label{wc:minimum}
 \sum_{0<Nu\le H}|E_u|^2\ll_{\nu,S,\delta,\varepsilon}
 \mathcal M_4(W)^2(HD)^\varepsilon
 \min\{D^{2\delta}U^{2-4\delta}H,\ H+DH^{1/3}+(HD)^{2/3}\}.
\end{equation}
At $H=D^{3/4}$ and $U=D^a$, $0\le a\le1/2$, the excess over $H$
is at most $\min\{2\delta+(2-4\delta)a,1/2\}$, up to arbitrarily
small losses. The new envelope improves the old one for $a>1/4$ as
$\delta\downarrow0$, but still leaves an exponent loss of $1/2$.
It does not settle the original signed short-row bound, type-II covariance,
positive arithmetic determinant comparison or GRH. The principal-residue
bound used here is sufficient for this envelope, not the stronger target
needed to reassemble the full signed argument.

\section{High cutoffs, actual bilinear factors and a signed covariance test}\label{sb:core}
Retain the original Eisenstein family, fixed finite-order $\nu$, fixed $S$,
common $W\in C_c^\infty((0,\infty))$ supported in $[1,2]$, and every raw
shared-prime zero. Put
\[
 C_U(m)=\sum_{\substack{abc=m\\Na>U,\,Nb>U}}\mu(a)\mu(b),\qquad
 r_U(m)=\#\{(a,b):ab=m,\ Na,Nb>U\}.
\]
All factors are ideals, squarefree and pairwise coprime; on the active
annulus $A_U=C_U-\mu$. A coefficient identity is different from a row-energy bound.

\begin{proposition}[The high-cutoff window]\label{sb:window}
If $D>1$ and $2D/7\le U^2\le D$, then for every surviving squarefree
column $D\le Nm\le2D$,
\[
 C_U(m)=\mu(m)r_U(m),\qquad
 A_U(m)=\mu(m)(r_U(m)-1),\qquad |A_U(m)|\ge1.       
\]
For each fixed nonzero $W$, the normalized squared coefficient norm satisfies
\begin{equation}
 D^{-1}\sum_{\substack{m\ {\rm sf}\\(m,S)=1}}
 |A_U(m)W(Nm/D)|^2\gg_{W,S}1                     \label{sb:coefficient-lower}
\end{equation}
for large $D$, uniformly in this window. This is not a lower bound for row energy.
\end{proposition}
\begin{proof}
The least norm of a nonunit ideal outside $S$ is at least $7$. A long-long
term with $c\ne1$ would have $Nm>7U^2\ge2D$. Thus $ab=m$, including both
endpoints. The fixed-point-free swap $(a,b)\mapsto(b,a)$ makes $r_U$ even,
so $A_U$ is an odd nonzero integer. The last claim follows by counting
squarefree ideals on a fixed interval where $|W|$ is bounded below; the
constant and error are proved in the extended proof.
\end{proof}

The exact remaining polynomial is genuinely bilinear, with $(a,b)=1$,
$U<Na,Nb<2D/U$, $D\le Na\,Nb\le2D$, and
$x_{ab}(u)=x_a(u)x_b(u)$, including zeros. Both factors see the same row
$u=\varepsilon\alpha\beta^2$. Cauchy followed by the accepted powered BGL
bound gives only
\begin{equation}
 Q_{B_X}(C_U,C_U)\ll_{\nu,S,\epsilon}
 (DX)^\epsilon\|W\|_\infty^2
 \{\sqrt D\,X^2+DX+D^{5/6}X^{5/3}\}.              \label{sb:factor-bound}
\end{equation}
At $X\asymp D^{1/4}$ the exponents are $1,5/4,5/4$; the earlier regrouped
bound has $1/2,5/4,13/12$ and is stronger. No $BT_\sigma$ saving is obtained.
Multiplying two marginal operator norms does not supply a coupled fourth moment.

\paragraph{A finite test in the original family.}
There is no automatic nonnegative signed covariance, even when
$C_U(m)\mu(m)\ge0$. Take $\nu=1$, $S$ the primes above $2,3$, $D=400$,
$U=17$, and all $72$ literal rows of $B_7$. Choose the common nonnegative
smooth profile constructed in the extended proof, with values $10$ at $439/400$, $1$ at
$475/400$, and zero at every other integer norm divided by $400$.
There are exactly four active squarefree ideal columns, of norms
$439,439,475,475$. In that order,
\[
 A=(1,1,1,1),\quad C=(0,0,2,2),\quad
 \mu=(-1,-1,1,1),\qquad
 \Re G_{B_7}=
 \begin{pmatrix}72&0&0&18\\0&72&18&0\\0&18&72&0\\18&0&0&72\end{pmatrix}.
\]
Exact raw-symbol enumeration, independently reproducible from the ideal
and residue-field definitions, gives
\begin{equation}
 Q(C,\mu)=-\frac{27}{25},\qquad
 Q(C,C)=\frac{36}{25},\quad
 Q(\mu,\mu)=\frac{864}{25},\quad
 Q(A,A)=\frac{954}{25}.                          \label{sb:covariance-example}
\end{equation}
The real Gram matrix is positive definite, with eigenvalues $54,90$;
its positivity does not force a cross covariance to be positive.
This fixed finite example refutes that shortcut, not the asymptotic
$BT_\sigma$ target or the original signed mean-square estimate.

\paragraph{The two-cutoff contour remains coupled.}
For $T_0=S\cup\operatorname{supp}(u)$, $\psi=\psi_u$ and independently
primitive $\psi_2=\operatorname{prim}(\psi^2)$, put $q_p=\psi(p)(Np)^{-s}$,
$\alpha_p=(Np)^{-z}$, $\beta_p=(Np)^{-w}$ outside $T_0$. The actual
three-placement Euler factor is $1+q_p-q_p\alpha_p-q_p\beta_p$. Its exact
factorization is
\begin{equation}
 \mathcal F(s,z,w)=
 \frac{L^{T_0}(s,\psi)}{L^{T_0}(s+z,\psi)L^{T_0}(s+w,\psi)}
 \prod_{p\notin T_0}\left(1-
 \frac{q_p^2(1-\alpha_p)(1-\beta_p)}
 {(1-q_p\alpha_p)(1-q_p\beta_p)}\right).           \label{sb:perron-factor}
\end{equation}
Exact double Perron inversion is valid on positive auxiliary lines
$c,d>1-\Re s$ for $1/2<\Re s<1$, retaining the original principal pole
exclusion. These safe lines reproduce the previous $U^{2-4\delta}$
energy loss at $\Re s=1/2+\delta$, with the same common-profile derivative
and fixed-margin dependencies. The correction product alone does not
justify moving those lines through reciprocal-$L$ poles. In fact
$\mathcal F(s,z,0)=1/L^{T_0}(s+z,\psi)$ as a meromorphic identity.
No zero-free region, reciprocal moment, analytic saving, new Lean
formalization or GRH proof is supplied by this construction.

\section{A quantitative covariance budget without a new exponent}\label{cbq:core}
Retain the original Eisenstein sextic symbols, fixed finite-order $\nu$,
fixed $S$, one possibly complex common $W\in C_c^\infty$ supported in
$[1,2]$, all squarefree columns prime to $S$, and every raw shared-prime
zero. In the accepted window $2D/7\le U^2\le D$, put $r_m=r_U(m)$,
$\lambda_m=\mu(m)\nu(m)W(Nm/D)$ and $v_m=\overline{\lambda_m}$.
Thus $C_U=\mu r$ and $A_U=\mu(r-1)$ on the active annulus.
For a finite literal row block $B$, let
\[
 G_{mn}=\sum_{u\in B}x_m(u)\overline{x_n(u)},\qquad
 R=\operatorname{diag}(r_m),\qquad
 M_j=\sum_mr_m^j|W(Nm/D)|^2\quad(0\le j\le2).
\]
The primitive $\psi_u,\psi_{u,2}$ and full deleted sets retain their
previous meanings; no conductor-drop value replaces a raw zero.

\begin{proposition}[An actual-vector lower budget]\label{cbq:budget}
If $0\le G\le\Lambda_B I$, then
\begin{equation}\label{cbq:vector}
 \Re Q_B(C_U,\mu)\ge
 -\frac{\Lambda_B}{2D}\{\sqrt{M_0M_2}-M_1\}.
\end{equation}
In this high window, uniformly in $U$,
\begin{equation}\label{cbq:moments}
 \sum_{D\le Nm\le2D}r_m\ll_K D,\qquad
 \sum_{D\le Nm\le2D}r_m^2\ll_K D\log(2D).
\end{equation}
Consequently
\[
 \Re Q_B(C_U,\mu)\ge
 -C_K\|W\|_\infty^2\Lambda_B\sqrt{\log(2D)}.
\]
\end{proposition}
The vector inequality is sharp among positive contractions with the
specified $v,Rv$; no such sharpness is asserted for the arithmetic Gram.
The moment bound counts actual balanced ideal factorizations, rather than
assuming independence of the two row factors.

On $B_X$, inserting the previously accepted powered-BGL norm gives only
\[
 \Re Q_{B_X}(C_U,\mu)\ge
 -C_{\nu,S,\epsilon}\|W\|_\infty^2(DX)^\epsilon
 \{X^2+DX+D^{2/3}X^{5/3}\}.
\]
At $X\asymp D^{1/4}$ the exponents remain $1/2,5/4,13/12$.
At $H=D^{3/4}$ this budget still has the old $D^{1/2}$ excess over $H$;
it supplies no $BT_\sigma$ saving. A negative-part budget is weaker than
$2\Re Q(C_U,\mu)\ge Q(C_U,C_U)+Q(\mu,\mu)-R_\sigma$.

For $S_R=\sqrt R$, the exact identity
\[
 \tfrac12(RG+GR)=S_RGS_R+\tfrac12[S_R,[S_R,G]]
\]
localizes the commutator to unequal factor-count levels, with entry
$(\sqrt{r_m}-\sqrt{r_n})^2G_{mn}$. If $B$ retains all six literal units,
the actual unit sum further gives $G_{mn}=0$ unless $Nm\equiv Nn\pmod{36}$.
These are exact deletions, not a power saving. Even zero different-level
pairings leave the nonnegative within-level energies
$\sum_k(k-1)^2Q_B(\mu\mathbf1_{r=k},\mu\mathbf1_{r=k})$.
The unequal-product coupled fourth moment and its mixed M\"obius term
remain unestimated. No new analytic exponent, energy lower bound,
Lean formalization or GRH conclusion follows.

\paragraph{Classical reference}NIST Digital Library of Mathematical Functions, Section25.15, \href{https://dlmf.nist.gov/25.15#E5}{Dirichlet L-functions, equations 25.15.5--6}; primitive functional equation and finite Gauss sums (snapshot pinned 4 November 2026). The normalization and explicit finite Gauss sums used here are stated above.\end{document}
